\documentclass[12pt]{article}

\input{../../_assets/preambulo.tex}


\begin{document}

    % 1. Foto de fondo
    % 2. Título
    % 3. Encabezado Izquierdo
    % 4. Color de fondo
    % 5. Coord x del titulo
    % 6. Coord y del titulo
    % 7. Fecha

    
    \input{../../_assets/portada}
     \portadaExamen{ffccA4.jpg}{Topología I\\Examen XIII}{Topología I. Examen XIII}{MidnightBlue}{-8}{28}{2026}{Roxana Acedo Parra}

    \begin{description}
        \item[Asignatura] Topología I.
        \item[Curso Académico] 2025-26.
        \item[Grado] Doble Grado en Ingeniería Informática y Matemáticas.
        \item[Grupo] Único.
        \item[Profesor] Antonio Alarcón López.
        \item[Descripción] Primer parcial.
        \item[Fecha] 30 de octubre de 2025.
        \item[Duración] 120 minutos.
    \end{description}
    
    \newpage
    
    \begin{ejercicio}[1 punto] Define los conceptos de punto adherente y adherencia de un conjunto $A$ en un espacio topológico $(X,\mathcal{T})$. Demuestra las siguientes propiedades:
        \begin{enumerate}[label=(\alph*)]
            \item La adherencia $\overline{A}$ de $A$ es cerrado en $(X,\mathcal{T})$.
            \item Si $C\subset X$ cerrado en $(X,\mathcal{T})$ y $A\subset C$ entonces $\overline{A}\subset C$
            \item $A$ es cerrado en $(X,\mathcal{T})$ si y solo si $A=\overline{A}.$
        \end{enumerate}
    \end{ejercicio}
    
    \begin{ejercicio}[2 puntos] Denotamos por $\mathcal{T}_{u}$ la topología usual en el conjunto $\mathbb{R}$ de los números reales.
        Dado un subconjunto $A\subset\mathbb{R}$ con $\emptyset\ne A\ne\mathbb{R},$ consideramos la familia de subconjuntos
        
        $$\mathcal{T}(A)=\{U\in\mathcal{T}_{u}:U\cap A=\emptyset\}\cup\{U\in\mathcal{T}_{u}:A\subset U\}\subset\mathcal{P}(\mathbb{R})$$
        
        Responde a 4 de los siguientes apartados:
        \begin{enumerate}[label=(\alph*)]
            \item Demuestra que $\mathcal{T}(A)$ es una topología en $\mathbb{R}$. Determina qué topología de entre $\mathcal{T}(A)$ y $\mathcal{T}_{u}$ es más fina.
            \item Calcula el interior y la adherencia de $[0, 1]$, $(0,\sqrt{2})$ y $(1,+\infty)$ en $(\mathbb{R},\mathcal{T}(\mathbb{N}))$ dónde $\mathbb{N}=\{1,2,3,...\}$
            \item Demuestra que el espacio topológico $(\mathbb{R},\mathcal{T}(\{0,1\}))$ es 1AN.
            \item Decimos que un espacio topológico $(X, \mathcal{T})$ es $T0$ (o de Kolmogórov) si para cada par de puntos $x, y\in X$ con $x\ne y$ se tiene que $\mathcal{N}_{x}\ne\mathcal{N}_{y},$ es decir, existe un entorno de uno de los puntos que no es entorno del otro.
            Demuestra que $(\mathbb{R},\mathcal{T}(A))$ es $T0$ si y solo si $A$ consiste de un único punto.
            \item Determina si el espacio topológico $(\mathbb{R}\setminus\mathbb{Q},\mathcal{T}(\mathbb{Q})_{|\mathbb{R}\setminus\mathbb{Q}})$ es $T1$.
        \end{enumerate}
    \end{ejercicio}

\newpage
\setcounter{ejercicio}{0}

    \begin{ejercicio}[1 punto] Define los conceptos de punto adherente y adherencia de un conjunto $A$ en un espacio topológico $(X,\mathcal{T})$. Demuestra las siguientes propiedades:
        \begin{enumerate}[label=(\alph*)]
            \item La adherencia $\overline{A}$ de $A$ es cerrado en $(X,\mathcal{T})$.
            \item Si $C\subset X$ cerrado en $(X,\mathcal{T})$ y $A\subset C$ entonces $\overline{A}\subset C$
            \item $A$ es cerrado en $(X,\mathcal{T})$ si y solo si $A=\overline{A}.$
        \end{enumerate}
        Teoría dada en clase.
    \end{ejercicio}
    
    \begin{ejercicio}[2 puntos] Denotamos por $\mathcal{T}_{u}$ la topología usual en el conjunto $\mathbb{R}$ de los números reales.
        Dado un subconjunto $A\subset\mathbb{R}$ con $\emptyset\ne A\ne\mathbb{R},$ consideramos la familia de subconjuntos
        
        $$\mathcal{T}(A)=\{U\in\mathcal{T}_{u}:U\cap A=\emptyset\}\cup\{U\in\mathcal{T}_{u}:A\subset U\}\subset\mathcal{P}(\mathbb{R})$$
        
        Responde a 4 de los siguientes apartados:
        \begin{enumerate}[label=(\alph*)]
            \item Demuestra que $\mathcal{T}(A)$ es una topología en $\mathbb{R}$. Determina qué topología de entre $\mathcal{T}(A)$ y $\mathcal{T}_{u}$ es más fina.
                Tenemos $\emptyset\ne A\subset \mathbb{R}$, $A\ne \mathbb{R}$, y
                
                $$\mathcal{T}(A) = \{U \in \mathcal{T}_u : U \cap A = \emptyset\} \cup \{U \in \mathcal{T}_u : A \subset U\}$$
    
                \textbf{¿$\mathcal{T}(A)$ topología en $\mathbb{R}$?} Comprobamos (A1), (A2), (A3)
                
                \begin{itemize}[leftmargin=*]
                    \item[\textbf{(A1)}] $\emptyset \in \mathcal{T}_u$, $\emptyset \cap A = \emptyset \implies \emptyset \in \mathcal{T}(A)$.\\
                    $\mathbb{R} \in \mathcal{T}_u$, $A \subset \mathbb{R} \implies \mathbb{R} \in \mathcal{T}(A)$.
                
                    \item[\textbf{(A2)}] Sea $\{U_i\}_{i \in I} \subset \mathcal{T}(A)$. ¿$\bigcup_{i \in I} U_i \in \mathcal{T}(A)$?\\
                    Como $\mathcal{T}_u$ es topología y $U_i \in \mathcal{T}_u$ $\forall i \in I$, tenemos que $\bigcup_{i \in I} U_i \in \mathcal{T}_u$.\\
                    Ahora si $U_i \cap A = \emptyset$ $\forall i \in I$, entonces:
                    
                    $$\left(\bigcup_{i \in I} U_i\right) \cap A = \bigcup_{i \in I} (U_i \cap A) = \emptyset \implies \bigcup_{i \in I} U_i \in \mathcal{T}(A)$$
                    
                    Si por el contrario $\exists j \in I$ tal que $A \subset U_j$, entonces $A \subset \bigcup_{i \in I} U_i$, luego $\bigcup_{i \in I} U_i \in \mathcal{T}(A)$.
                
                    \item[\textbf{(A3)}] Sean $U_1, U_2 \in \mathcal{T}(A)$. ¿$U_1 \cap U_2 \in \mathcal{T}(A)$?\\
                    Como $\mathcal{T}_u$ es topología y $U_1, U_2 \in \mathcal{T}_u$, tenemos que $U_1 \cap U_2 \in \mathcal{T}_u$.\\
                    Ahora si $A \subset U_1$ y $A \subset U_2$, entonces $A \subset U_1 \cap U_2$, luego $U_1 \cap U_2 \in \mathcal{T}(A)$.\\
                    Por el contrario $\exists i \in \{1, 2\}$ tal que $A \cap U_i = \emptyset$, entonces $A \cap (U_1 \cap U_2) = \emptyset$, luego $U_1 \cap U_2 \in \mathcal{T}(A)$.
                \end{itemize}
                
                Como $\mathcal{T}(A)$ cumple (A1), (A2) y (A3), es una topología.
                
                \textbf{¿Qué topología de entre $\mathcal{T}(A)$ y $\mathcal{T}_u$ es más fina?}
                
                Por definición de $\mathcal{T}(A)$, todos los abiertos de $\mathcal{T}(A)$ lo son de $\mathcal{T}_u$ (es decir, $U \in \mathcal{T}(A) \implies U \in \mathcal{T}_u$), luego $\mathcal{T}(A) \le \mathcal{T}_u$.\\
                En general, $\mathcal{T}_u \nsubseteq \mathcal{T}(A)$, esto solo ocurre cuando $A$ solo tiene un elemento, en ese caso $\mathcal{T}(A) = \mathcal{T}_u$.
            
            \item Calcula el interior y la adherencia de $[0, 1]$, $(0,\sqrt{2})$ y $(1,+\infty)$ en $(\mathbb{R},\mathcal{T}(\mathbb{N}))$ dónde $\mathbb{N}=\{1,2,3,...\}$
            
                $$\mathcal{T}(\mathbb{N}) = \{U \in \mathcal{T}_u : U \cap \mathbb{N} = \emptyset\} \cup \{U \in \mathcal{T}_u : \mathbb{N} \subset U\}$$

                \begin{itemize}[leftmargin=*]
                    \item \textbf{De $[0, 1]$:}
                    
                    $$[0, 1]^\circ = (0, 1)$$
                    
                    En efecto: $(0, 1) \in \mathcal{T}_u$ y $(0, 1) \cap \mathbb{N} = \emptyset \implies (0, 1) \in \mathcal{T}(\mathbb{N})$.\\
                    Además $0$ y $1$ no son puntos interiores de $[0, 1]$ ya que cualquier entorno de uno de estos puntos contiene un abierto de $\mathcal{T}_u$, luego un intervalo abierto, que contiene al punto, y un intervalo así no puede estar contenido en $[0, 1]$.\\
                    En conclusión, $(0, 1)$ es el mayor abierto en $\mathcal{T}(\mathbb{N})$ dentro de $[0, 1]$.
                
                    
                    $$\overline{[0, 1]} = [0, 1] \cup \mathbb{N}$$
                    
                    En efecto: $[0, 1] \cup \mathbb{N} \in \mathcal{C}_{\mathcal{T}(\mathbb{N})}$ porque su complementario está en $\mathcal{T}_u$ y no corta a $\mathbb{N}$.\\
                    Además, si $x \in \mathbb{N}$ y $U \in \mathcal{T}(\mathbb{N})$ con $x \in U$, entonces $U \notin \mathcal{T}_u \implies \mathbb{N} \subset U$, luego $U$ es entorno de $1$. Esto implica que $x \in \overline{[0, 1]}$.\\
                    En conclusión, $[0, 1] \cup \mathbb{N}$ es el cerrado más pequeño que contiene a $[0, 1]$.
                
                    \item \textbf{Usando argumentos similares a los anteriores, tenemos que:}
                    
                    $$(0, \sqrt{2})^\circ = (0, \sqrt{2}) \setminus \{1\}$$
                    
                    (En efecto, $(0, \sqrt{2}) \setminus \{1\} \in \mathcal{T}(\mathbb{N})$ porque está en $\mathcal{T}_u$ y no corta a $\mathbb{N}$, y $1$ no es interior a $(0, \sqrt{2})$ porque cualquier entorno suyo contiene a $\mathbb{N}$.)
                    
                    $$\overline{(0, \sqrt{2})} = [0, \sqrt{2}] \cup \mathbb{N}$$
                    
                    $$(1, +\infty)^\circ = (1, +\infty) \setminus \mathbb{N}$$
                    
                    $$\overline{(1, +\infty)} = [1, +\infty)$$
                \end{itemize}

            \item Demuestra que el espacio topológico $(\mathbb{R},\mathcal{T}(\{0,1\}))$ es 1AN.

                $$\mathcal{T}(\{0, 1\}) = \{U \in \mathcal{T}_u : U \cap \{0, 1\} = \emptyset\} \cup \{U \in \mathcal{T}_u : \{0, 1\} \subset U\} $$
                Tenemos que encontrar de cada punto $x \in \mathbb{R}$ una base de entornos numerable:
                
                \begin{itemize}[leftmargin=*]
                    \item Si $x \notin \{0, 1\}$, tomemos por ejemplo:
                    
                    $$\beta_x = \left\{ \left(x - \frac{1}{n}, x + \frac{1}{n}\right) : n \in \mathbb{N} \right\}$$
                    
                    ¿Es $\beta_x$ b.d.e. de $x$?\\
                    $\left(x - \frac{1}{n}, x + \frac{1}{n}\right)$ contiene en $x$ un intervalo centrado y disjunto de $\{0, 1\}$, luego $\left(x - \frac{1}{n}, x + \frac{1}{n}\right)$ es entorno de $x$ en $\mathcal{T}(\{0, 1\})$. Así, $\beta_x \subset \mathcal{N}_x$.\\
                    Ahora; si $N \in \mathcal{N}_x$, $\exists U \in \mathcal{T}(\{0, 1\})$ con $x \in U \subset N$.\\
                    Como $U \in \mathcal{T}_u$, $\exists n \in \mathbb{N}$ con $\left(x - \frac{1}{n}, x + \frac{1}{n}\right) \subset U \subset N \implies \beta_x$ b.d.e.
                
                    \item Si $x \in \{0, 1\}$, tomamos por ejemplo:
                    
                    $$\beta_x = \left\{ \left(-\frac{1}{n}, \frac{1}{n}\right) \cup \left(1 - \frac{1}{n}, 1 + \frac{1}{n}\right) : n \in \mathbb{N} \right\}$$
                    
                    ¿Es $\beta_x$ b.d.e. de $x$?\\
                    Primero, cada elemento de $\beta_x$ contiene a $\{0, 1\}$, es abierto en $\mathcal{T}_u$ y luego abierto en $\mathcal{T}(\{0, 1\})$. Así, $\beta_x \subset \mathcal{N}_x$.\\
                    Ahora si $N \in \mathcal{N}_x$, con $x \in U \subset N$, tenemos que no puede ser $U \cap \{0, 1\} = \emptyset$, así que se cumple que $\{0, 1\} \subset U$. Además, $U \in \mathcal{T}_u$, así que $\exists n \in \mathbb{N}$ con $\left(-\frac{1}{n}, \frac{1}{n}\right) \cup \left(1 - \frac{1}{n}, 1 + \frac{1}{n}\right) \subset U \subset N \implies \beta_x$ b.d.e. de $x$.
                \end{itemize}
                
                Como $\beta_x$ es numerable $\forall x \in \mathbb{R}$, tenemos que $(\mathbb{R}, \mathcal{T}(\{0, 1\}))$ es 1AN.
                
            \item Decimos que un espacio topológico $(X, \mathcal{T})$ es $T0$ (o de Kolmogórov) si para cada par de puntos $x, y\in X$ con $x\ne y$ se tiene que $\mathcal{N}_{x}\ne\mathcal{N}_{y},$ es decir, existe un entorno de uno de los puntos que no es entorno del otro.
            Demuestra que $(\mathbb{R},\mathcal{T}(A))$ es $T0$ si y solo si $A$ consiste de un único punto.
            
                ¿$(\mathbb{R}, \mathcal{T}(A))$ es $T0$ para algún $A$? ¿o para algún $A$?
                \begin{itemize}[leftmargin=*]
                    \item[$\implies$] Sup. $\exists x, y \in A$ son $x \ne y$.\\
                    Sea $N \in \mathcal{N}_x$. Entonces $\exists U \in \mathcal{T}(A)$ con $x \in U \subset N$.\\
                    Como $x \in A$ y $x \in U$, tenemos que $A \subset U$, luego $y \in U \subset N$.\\
                    Esto implica que $N \in \mathcal{N}_y$. Así, $\mathcal{N}_x \subset \mathcal{N}_y$.\\
                    El mismo razonamiento demuestra que $\mathcal{N}_y \subset \mathcal{N}_x$, luego $\mathcal{N}_x = \mathcal{N}_y$.\\
                    Por tanto $(\mathbb{R}, \mathcal{T}(A))$ no es $T0$.
                
                    \item[$\impliedby$] Si $A = \{x_0\}$, entonces $\mathcal{T}(A) = \{U \in \mathcal{T}_u : x_0 \notin U\} \cup \{U \in \mathcal{T}_u : x_0 \in U\} = \mathcal{T}_u$.\\
                    Pero $(\mathbb{R}, \mathcal{T}_u)$ es claramente $T0$, de hecho es $T_2$, luego $(\mathbb{R}, \mathcal{T}(A)) = (\mathbb{R}, \mathcal{T}_u)$ es $T0$.
                \end{itemize}
            
            \item Determina si el espacio topológico $(\mathbb{R}\setminus\mathbb{Q},\mathcal{T}(\mathbb{Q})_{|\mathbb{R}\setminus\mathbb{Q}})$ es $T1$.

                ¿$(\mathbb{R}\setminus\mathbb{Q}, \mathcal{T}(\mathbb{Q})_{|\mathbb{R}\setminus\mathbb{Q}})$ es $T1$?

                Un e.t. es $T1$ si y solo si los conjuntos con un único elemento son cerrados.\\
                Para demostrar que $(\mathbb{R}\setminus\mathbb{Q}, \mathcal{T}(\mathbb{Q})_{|\mathbb{R}\setminus\mathbb{Q}})$ es $T1$, basta comprobar que $\forall x \in \mathbb{R}\setminus\mathbb{Q}$, $\{x\}$ es cerrado en ese espacio. Su complementario en $\mathbb{R}\setminus\mathbb{Q}$ es $(\mathbb{R}\setminus\mathbb{Q}) \setminus \{x\}$.\\
                Ahora, como $\mathbb{R} \setminus \{x\} \in \mathcal{T}_u$ y $\mathbb{Q} \subset \mathbb{R} \setminus \{x\}$, tenemos que $\mathbb{R} \setminus \{x\} \in \mathcal{T}(\mathbb{Q})$, luego $(\mathbb{R}\setminus\mathbb{Q}) \setminus \{x\}$ es la intersección de $\mathbb{R}\setminus\mathbb{Q}$ con un abierto de $\mathcal{T}(\mathbb{Q})$, esto es, $(\mathbb{R}\setminus\mathbb{Q}) \setminus \{x\} \in \mathcal{T}(\mathbb{Q})_{|\mathbb{R}\setminus\mathbb{Q}}$.\\
                Por tanto $\{x\} \in \mathcal{C}_{\mathcal{T}(\mathbb{Q})_{|\mathbb{R}\setminus\mathbb{Q}}}$, lo que demuestra que el espacio es $T1$.
        \end{enumerate}
    \end{ejercicio}
\end{document}